% Generated draft: requires mathematical and source review.
\documentclass[11pt]{article}
\usepackage[margin=25mm]{geometry}
\usepackage{fontspec}
\setmainfont{Latin Modern Roman}
\usepackage{amsmath,amssymb,mathtools,unicode-math}
\setmathfont{Latin Modern Math}
\usepackage{enumitem,needspace,xurl,hyperref}
\setlength{\parindent}{0pt}
\setlength{\parskip}{5pt}
\setlength{\emergencystretch}{3em}
\newcommand{\N}{\mathbb{N}}
\newcommand{\Z}{\mathbb{Z}}
\newcommand{\Q}{\mathbb{Q}}
\newcommand{\R}{\mathbb{R}}
\newcommand{\C}{\mathbb{C}}
\newcommand{\F}{\mathbb{F}}
\newcommand{\A}{\mathbb{A}}
\newcommand{\PP}{\mathbb{P}}
\newcommand{\E}{\mathbb{E}}
\newcommand{\eps}{\varepsilon}
\newcommand{\dd}{\,\mathrm{d}}
\begin{document}
\section*{UnsolvedMath ID: 30000856}
\textbf{Source problem code:} catalogue entry
\subsection{Definitions and mathematical statement}
Let \(X\) be a smooth complex projective variety of dimension \(n\geq1\), with cotangent bundle \(\Omega_X^1\), tangent bundle \(T_X\), and canonical line bundle \(K_X=\det\Omega_X^1\). A line bundle is nef when its degree on every integral curve is nonnegative; a vector bundle \(E\) is nef when the tautological line bundle on its projective bundle of one-dimensional quotients is nef. For a possibly singular curve \(C\), the usual restriction convention tests the pullback of \(E|_C\) to its normalization.

A covering family \((C_t)_{t\in T}\) is a family of curves whose universal-family evaluation map to \(X\) is dominant. The source additionally requires this family to be ``maximal,'' without defining that qualification. The first question is whether
\[
 \bigl[\exists\text{ a maximal covering family }(C_t)
     \text{ with }\Omega_X^1|_{C_t}\text{ non-nef for general }t\bigr]
 \ \Longrightarrow\ X\text{ is uniruled}.
\]
Uniruled means that a general point of \(X\) lies on a rational curve; equivalently in the usual projective setting there is a dominant rational map \(\PP^1\times Y\dashrightarrow X\) with \(\dim Y=n-1\). The family in this question is not silently replaced by a complete-intersection family.

For the second question assume \(-K_X\) is nef. Given an ample polarization tuple \(H=(H_1,\ldots,H_{n-1})\), an MR-general curve is a sufficiently general complete intersection
\[
 C=D_1\cap\cdots\cap D_{n-1},\qquad
 D_i\in|m_iH_i|,
\]
with the integers \(m_i\) sufficiently large in the intended Mehta--Ramanathan sense. For \(n=1\), take \(C=X\). The bundle \(T_X\) is generically nef with respect to \(H\) if its restriction to such a general curve is nef. The source asks both the weaker and stronger polarization questions:
\[
 -K_X\text{ nef}\quad\Longrightarrow\quad
 \exists H\text{ with }T_X\text{ generically nef with respect to }H\ ?
\]
\[
 -K_X\text{ nef}\quad\Longrightarrow\quad
 \forall H,\ T_X\text{ generically nef with respect to }H\ ?
\]
The record does not decide whether polarization means one ample class repeated or arbitrary ample tuples, or give all the size/genericity conventions. These questions concern nefness on MR-general curves, not nefness of \(T_X\) on every curve. The anticanonical hypothesis belongs to the second question and is not imposed on the first.
\subsection{Short English statement}
The first question asks whether non-nef cotangent restrictions along a specified covering family force rational curves through general points. The second asks whether nef anticanonical class forces generic tangent nefness for at least one, or for every, polarization. The family maximality and polarization conventions remain incomplete.
\subsection{Sources}
\begin{itemize}
\item Local UnsolvedMath registry, permanent ID \(30000856\), exact source code \texttt{catalogue entry}, proposed year 2007; no individual proposer is supplied.
\item \emph{Komplexe Algebraische Geometrie (2008)}, Oberwolfach Reports, DOI 10.4171/owr/2007/47, pp. 2791--2850. \url{https://doi.org/10.4171/owr/2007/47}.
\item Supplied reference: Ou, ``generic nefness of tangent sheaves,'' \url{https://arxiv.org/abs/1703.03175}.
\item The supplied \texttt{corrected\_verified} note reports checking mathematical content against the official report. The assessment checked 2026-08-21 describes a combined fragment missing hypotheses and says only partial surrounding results and no complete resolution were verified for it. The open label and assessment are dated supplied triage requiring independent review, including comparison with the supplied clean statement.
\end{itemize}
\end{document}
